\section{Remaining Steps}
\label{sec:remaining_steps}

\paragraph{Baseline testing}
We have set up our workspace. Our repository is a fork of Chipyard~1.15.0 with
SonicBOOM included in the tree. A single wrapper script sets up the Chipyard toolchain,
Verilator, Spike and Vivado~2023.2 on the department's Linux servers, so every build,
simulation and synthesis run uses the same environment. We have also tagged the
unmodified Chipyard~1.15.0 with SonicBOOM as our frozen baseline, and our next step is
to measure it. We will build the baseline in Verilator and check that it passes the
RISC-V ISA tests and the riscv-tests benchmarks. We will then port the BOOM Spectre
proofs of concept to Chipyard bare-metal tests and show that each one recovers the
secret on the baseline, which shows that the tests themselves work. At the same time,
we will record the baseline cycles and IPC on CoreMark, Dhrystone and the riscv-tests
benchmarks, and run the Vivado flow to get its utilization, timing slack and power
estimates. These results give us a reference point for all three axes of our
evaluation, which are security, performance and hardware cost. Every later number will
be compared to this baseline and reported with the commit it was measured on.

\paragraph{Development}
Development will follow the order of Table~\ref{tab:configs}. The conservative
stand-in comes first. Track B can build it without waiting for track A, and it lets
security testing begin early and gives an upper bound on the performance cost.
Meanwhile, track A builds the taint table at rename using the interface shared between
the two tracks (Section~\ref{sec:approach}). Baseline STT comes next, where track B
replaces the conservative stand-in's branch-mask condition with the
\textit{isSafe(YRoT)} check and makes every load/store unit path that can fire a load
respect taint. The two $\mu$STT mechanisms come last. The released $\mu$STT
source~\cite{Chen:uSTT:2025} is included in our repository as a read-only reference.
We will port the Age Matrix to Chisel at BOOM's widths and check that it produces the
same YRoTs as the baseline STT bypass. For the impede micro-op, we will evaluate the
two options from Section~\ref{sec:approach}, which are emitting an extra micro-op at
dispatch or reserving a dedicated issue-queue entry. Each task is tracked as a GitHub
issue and developed on its own branch, then merged by pull request. Any change to a
file shared between tracks needs a review from another teammate before it is merged.
Each configuration must pass the checks in Section~\ref{sec:verification} before we
report its performance and hardware cost. Its results are then posted on its issue
along with the configuration and commit that produced them.
